\scanpage{661}
\mspage{49}
% archive sheet number: 660
cardinale etc).

\uline{Question 3.15} Cela me ramène à la question particulièrement \unsure{élusive} et irritante, de l’existence des petits ensembles $Q$-générateurs, dans une partie $\Phi$ de $\mathrm{Fl}\,M$, quand on suppose (disons) que $M$ est stable par petites $\varprojlim$, \inserted{que ses éléments sont accessibles} et \inserted{qu’il} admet une filtration cardinale.

Ainsi on a la question : une partie $L$-accessible de $\mathrm{Fl}\,M$, \inserted{\deleted{\unsure{q|p}}\unsure{disons -pré au besoin|disons -pour au besoin}} $Q$-saturée \deleted{\unsure{trans|traca}} à gauche (ou ce qui revient au $=$, contenant les iso, et stable par cochgt de base, par composition et par facteurs directs) admet-elle une \uline{petite} partie $\Phi_0$ $Q$-génératrice à g. ?

{\reversemarginpar\margin{!\hfill\rule[-3em]{.4pt}{6em}}}
D’autre part, je n’ai pas à présent construit un seul exemple d’un couple $\Phi\longleftrightarrow\Psi$, dans une catégorie $M$, qui ne soit factorisant ! Le cas où $M$ est spéciale, étant bien sûr le plus intéressant. Visiblement, dans tout
